\documentclass[11pt]{article}
\usepackage[a4paper,margin=28mm]{geometry}
\usepackage{amsmath,amssymb,amsthm,mathtools,hyperref}
\newcommand{\ddc}{dd^c}
\newcommand{\tr}{\operatorname{tr}}
\newcommand{\Ga}{\mathcal G}
\newtheorem{theorem}{Theorem}
\newtheorem{lemma}[theorem]{Lemma}
\theoremstyle{remark}\newtheorem{remark}[theorem]{Remark}
\title{The complete Gauduchon determinant profile\newline
and a limitation on pointwise ellipticity}
\author{}\date{}
\begin{document}\maketitle
\noindent\textbf{Working draft.} Independent mathematical review and
reconstruction are pending. The unrestricted balanced--pluriclosed
coexistence and flow questions of P13 remain unresolved.
\begin{abstract}
For a fixed smooth probability volume and a positive Hermitian background
whose lower powers are $dd^c$-closed,
the normalized Gauduchon-current determinant profile is unchanged by
a positive global scalar-potential perturbation of the metric.
The supremum over all positive currents, including singular and
degenerate ones, is computed exactly on every density root in $L^n$.
On a fixed complex torus we construct both K\"ahler metrics and genuinely
nonclosed pluriclosed metrics with identical complete determinant data
and positive constant bounded-potential mass, but with pointwise
ellipticity tending to zero. The nonclosed family has fixed nonzero
exterior derivative and uniformly bounded smooth coefficients.
Another family consists of sharp all-current determinant maximizers
with uniformly bounded smooth potentials and a continuous canonical
limit. Its determinant-power densities concentrate, while the limiting
current has zero intrinsic determinant measure and retains positive
bounded-potential Monge--Amp\`ere mass.
This specifies a limitation of the determinant data alone. It supplies
neither a non-K\"ahler coexistence example nor an assertion about any
trajectory of pluriclosed flow.
\end{abstract}

\section{The current family and its intrinsic determinant measure}
Use $\ddc=i\partial\bar\partial/(2\pi)$, ordinary wedge powers and
canonical psh-plus-smooth representatives. Fix a connected compact
complex manifold $X$ of dimension $n\geq2$ and a smooth positive
probability volume $dV$. For a smooth positive Hermitian metric $h$ put
\[
 \Ga_h=\{T\geq0\text{ of bidegree }(n-1,n-1):
            \ddc T=0,\ \int_Xh\wedge T=1\}.
\]
The family consists of actual currents, not only smooth forms.
For such a positive current write $\sigma_h=h\wedge T$. In an
$h$-unitary frame let $P_h$ be its trace-normalized matrix density,
so $\tr P_h=n$ and
\[
 \xi\wedge T=\frac1n\tr(BP_h)\,d\sigma_h
\]
for a smooth real $(1,1)$-form $\xi$ with $h$-relative matrix $B$.
The determinant measure is
\[
 \Gamma_{0,h}(T)=
 \left(\frac{dV}{h^n}\right)^{1/n}
       \det(P_h)^{1/n}\sigma_h.
\]
For now it is a Radon measure; its absolute continuity will be proved
under the differential hypothesis below.

\begin{lemma}\label{lem:intrinsic}
For a fixed positive current $T$ and fixed $dV$, the measure
$\Gamma_{0,h}(T)$ is independent of the positive auxiliary metric $h$.
This statement includes singular and degenerate currents.
\end{lemma}
\begin{proof}
In a holomorphic frame put
$\lambda=\bigwedge_{j=1}^n(i\,dz_j\wedge d\bar z_j)$,
$h=i\sum H_{ij}\,dz_i\wedge d\bar z_j$ and $dV=v\lambda$.
Positive currents have measure coefficients. Choose a positive
Radon measure $\nu$ dominating them and the positive Hermitian
coefficient matrix $Q$ characterized by
$\xi\wedge T=\tr(\Xi Q)\nu$. Then
\[
 \sigma_h=\tr(HQ)\nu,\qquad
 P_h=\frac{nH^{1/2}QH^{1/2}}{\tr(HQ)}.
\]
The formula holds on $Q\ne0$; zero is assigned when $Q=0$.
Since $h^n=n!\det(H)\lambda$, determinant algebra gives
\begin{equation}\label{eq:intrinsic}
 \Gamma_{0,h}(T)
 =n(n!)^{-1/n}v^{1/n}(\det Q)^{1/n}\nu.
\end{equation}
This expression has no $H$. Replacing $\nu$ by $L\nu$ replaces
$Q$ by $Q/L$, leaving the measure unchanged; a common dominating
measure compares any two choices. The original geometric definition
ensures compatibility under frame changes. Matrix AM--GM bounds it
locally by a constant times a reference trace measure, proving local
finiteness. Compactness gives a finite Radon measure.
\end{proof}
Write $\Gamma_0(T)$ for this intrinsic measure.

\section{Bounded masses and the exact entire-family endpoint}
\begin{lemma}\label{lem:mass}
Let $\beta$ be smooth and positive with
$\ddc(\beta^j)=0$ for $1\leq j\leq n-1$, and put
$V=\int_X\beta^n$. Every bounded canonical $\beta$-psh $w$ satisfies
\[
 \int_X(\beta+\ddc w)^n=V.
\]
For a smooth $w$ the lower powers of $A=\beta+\ddc w$ are also
$dd^c$-closed, and $\int_X\beta\wedge A^{n-1}=V$.
No $d$-closedness of $\beta$ is required.
\end{lemma}
\begin{proof}
Put $E=\ddc w$. Locally choose a smooth strictly psh correction $q$
so that $w+q$ is bounded psh. Expand the signed powers of
$E=\ddc(w+q)-\ddc q$ using bounded Bedford--Taylor mixed products
\cite[equation (2.1), Theorem 2.1 and Corollary 2.2]{BT}.
To see the recursion explicitly put $S=\ddc(w+q)$ and $R=\ddc q$.
Expand $(S-R)^{b-1}$; for each closed mixed factor $U$,
the bounded recursion gives $\ddc[(w+q)U]=S\wedge U$, and smooth
differentiation gives $\ddc(qU)=R\wedge U$. Subtraction and the
binomial identity give
\[
 E^b=\ddc\bigl[w E^{b-1}\bigr],\qquad b\geq1.
\]
The same smooth-correction identities show agreement on overlaps,
so these are global closed currents. The bounded canonical
representative specifies the multiplier.
For $a+b=n$, $b\geq1$, pairing with the smooth form $\beta^a$ yields
\[
 \int_X\beta^a\wedge E^b
   =\int_Xw\,\ddc(\beta^a)\wedge E^{b-1}=0.
\]
Closedness of $E^{b-1}$ removes the mixed derivative terms; for $a=0$
the derivative of the constant is zero. The genuine bounded Hermitian
product is the binomial expansion of these local products with the
smooth $\beta$ factors. Its total mass is therefore $V$.
For smooth $w$, expand $A^j$ and apply $dd^c$ termwise. All factors
containing $E$ are closed, so only $\ddc(\beta^{j-b})$ remains.
This vanishes for $j\leq n-1$. The same mixed-integral calculation
proves $\int\beta\wedge A^{n-1}=V$.
\end{proof}

\begin{theorem}\label{th:profile}
Let $\beta$ be smooth positive Hermitian with
$\ddc(\beta^j)=0$ for $1\leq j\leq n-1$,
$V=\int_X\beta^n$, and let $h=\beta+\ddc u>0$ for a smooth
global real function $u$. Then
\[
 \Ga_h=\Ga_\beta,\qquad
 \{\Gamma_{0,h}(T):T\in\Ga_h\}
 =\{\Gamma_{0,\beta}(T):T\in\Ga_\beta\}.
\]
Every measure in this common family has a density $k_T$ satisfying
\[
 k_T\geq0,\qquad
 \|k_T\|_{L^{n/(n-1)}(dV)}\leq V^{-1/n}.
\]
For every real or complex $p\in L^n(dV)$,
\begin{equation}\label{eq:profile}
 \sup_{T\in\Ga_h}\int_X|p|\,\Gamma_0(T)
       =V^{-1/n}\|p\|_{L^n(dV)}.
\end{equation}
All suprema retain singular and degenerate tests.
\end{theorem}
\begin{proof}
Distributional integration by parts gives
$\int h\wedge T=\int\beta\wedge T+\int u\ddc T
=\int\beta\wedge T$. This proves equality of the actual current
sets; Lemma~\ref{lem:intrinsic} proves equality of their measure sets.

Take a smooth strictly positive real $r$ and put $M=\int r^n dV$.
Smooth Hermitian Monge--Amp\`ere existence supplies a smooth
$A=\beta+\ddc w>0$ with
\begin{equation}\label{eq:smooth}
 A^n=\frac{Vr^n}{M}\,dV.
\end{equation}
Precisely, \cite[Corollary 1]{TW} applies on this connected compact
Hermitian manifold with
$F=\log(Vr^n/(M\beta^n/dV))$. The source potential is multiplied
by $2\pi$ for our operator. Lemma~\ref{lem:mass} gives source scalar
$b=0$, because both sides have total mass $V$. This step uses the
lower-power differential hypothesis, including when $\beta$ is nonclosed.
No uniform estimate on degenerating $F$, or flow theorem, is used.

For every $T\in\Ga_\beta$, integration by parts gives
$\int A\wedge T=1$. In an $A$-unitary frame, matrix AM--GM gives
$\det(P_A)^{1/n}\leq1$, including singular coefficient densities.
Thus as Radon measures
\[
 \Gamma_0(T)\leq
       (dV/A^n)^{1/n}(A\wedge T).
\]
Multiplying by $r$ and using \eqref{eq:smooth} yields
\begin{equation}\label{eq:dualbound}
 \int r\,\Gamma_0(T)\leq V^{-1/n}\|r\|_{L^n(dV)}.
\end{equation}
Replacing a smooth nonnegative $r$ by $r+\eta$ and letting
$\eta\downarrow0$ extends this bound to nonnegative smooth tests.
For smooth real or complex $s$, smooth positive approximations
$\sqrt{|s|^2+\eta^2}$ imply
$|\int s\,\Gamma_0(T)|\leq V^{-1/n}\|s\|_n$.
Consequently integration against $\Gamma_0(T)$ extends continuously
from smooth functions to $L^n$. Its positive dual representative is
$k_T\in L^{n/(n-1)}$ with the stated bound. Equality on smooth
tests identifies the two Radon measures, proving absolute continuity.
The integrals of arbitrary $L^n$ roots are therefore well-defined.

For the reverse inequality with $r>0$ smooth, let
$T_r=A^{n-1}/V$. It is smooth positive and $dd^c$-closed, and
Lemma~\ref{lem:mass} gives $\int\beta\wedge T_r=1$.
Its positive root is $\gamma=V^{-1/(n-1)}A$, so
\[
 k_{T_r}=(\gamma^n/dV)^{(n-1)/n}
       =V^{-1/n}\frac{r^{n-1}}{\|r\|_n^{n-1}}.
\]
It attains equality in \eqref{eq:dualbound}. For a nonnegative
$p\in L^n$, take smooth positive $r_j\to p$ in $L^n$. The common
dual bound controls $\int(p-r_j)k_{T_{r_j}}dV$, proving the lower
bound in \eqref{eq:profile}; the upper bound already holds for
every current. Apply the same argument to $|p|$ for complex roots.
\end{proof}

\section{Fixed complete data with degenerating ellipticity}
Here $v_-(h)$ and $v_+(h)$ denote the infimum and supremum of
$\int_X(h+\ddc w)^n$ over bounded canonical $h$-psh potentials $w$.
\begin{theorem}\label{th:degeneration}
On $X=(\mathbb C/(\mathbb Z+i\mathbb Z))^n$, $n\geq2$, there
are positive metrics $\omega_\varepsilon$ in a fixed K\"ahler
class, $0<\varepsilon\leq1$, with fixed complete current and
determinant-measure families and fixed endpoint \eqref{eq:profile},
while their least relative eigenvalue equals $\varepsilon$.
They have uniformly bounded smooth coefficients and global
potentials, and
\[
 v_-(\omega_\varepsilon)=v_+(\omega_\varepsilon)=n!>0.
\]
The manifold has a global holomorphic volume form and $c_1=0$.
Every displayed metric is already balanced and pluriclosed.
\end{theorem}
\begin{proof}
Write $z_j=x_j+iy_j$,
$\beta_j=(i/2)dz_j\wedge d\bar z_j=dx_j\wedge dy_j$,
$\beta=\sum_j\beta_j$, $V=n!$, and $dV=\beta^n/V$.
Set
\[
 u_\varepsilon=\frac{1-\varepsilon}{4\pi}\cos(4\pi x_1),
 \qquad a_\varepsilon=1-(1-\varepsilon)\cos(4\pi x_1).
\]
For a function of $x_1$, $\ddc u=(u''/(4\pi))\beta_1$.
Hence
\[
 \omega_\varepsilon=\beta+\ddc u_\varepsilon
 =a_\varepsilon\beta_1+\sum_{j=2}^n\beta_j.
\]
These are globally periodic smooth positive closed forms with
$\varepsilon\leq a_\varepsilon\leq2-\varepsilon$ and
$\omega_\varepsilon\leq2\beta$. Their minimum relative eigenvalue
is $\varepsilon$. They converge smoothly to a semipositive form
losing rank on $\{x_1\in\tfrac12\mathbb Z\}$.
All fixed-order smooth norms of the coefficients and potentials are
uniformly bounded. Theorem~\ref{th:profile} proves equality of the
entire current data and the exact norm at every $\varepsilon$.

For any bounded canonical $\omega_\varepsilon$-psh function $w$,
closed bounded Bedford--Taylor products \cite[Section 2]{BT} and their
defining recursion give
\[
 (\omega_\varepsilon+\ddc w)^j
 =\omega_\varepsilon\wedge(\omega_\varepsilon+\ddc w)^{j-1}
    +\ddc\bigl[w(\omega_\varepsilon+\ddc w)^{j-1}\bigr].
\]
Pairing with the remaining closed powers and applying distributional
Stokes removes each derivative term. The resulting total mass is
$\int\omega_\varepsilon^n=V$. The explicit formula gives the same
identity because the mean of $a_\varepsilon$ equals one.
Thus both extrema over bounded
canonical potentials equal $V$. The global form
$dz_1\wedge\cdots\wedge dz_n$ trivializes the canonical bundle;
closed positive metrics are balanced and pluriclosed.
\end{proof}

\section{The same obstruction for genuinely nonclosed metrics}
\begin{theorem}\label{th:nonclosed}
On the fixed torus above, fix $0<\tau<1/2$. There is a positive
nonclosed pluriclosed metric $h_\tau$ and a family
$\Omega_\varepsilon=h_\tau+\ddc u_\varepsilon>0$,
$0<\varepsilon\leq1/2$, with identical actual current sets and
intrinsic determinant-measure families. For any fixed smooth positive
probability volume their entire-family endpoint is
\[
 \sup_{T\in\Ga_{\Omega_\varepsilon}}
       \int_X|p|\,\Gamma_0(T)
       =V_\tau^{-1/n}\|p\|_{L^n(dV)},\qquad
 V_\tau=n!(1-\tau^2/2).
\]
Both bounded-potential masses equal $V_\tau$. All fixed-order smooth
coefficient and potential norms are uniformly bounded, while
\[
 \min_X\lambda_{\min}(\Omega_\varepsilon\mid\beta)=\varepsilon,
 \qquad d\Omega_\varepsilon=dh_\tau\ne0.
\]
The flat $\beta$ is a fixed balanced metric on this same manifold
with $c_1=0$; each $\Omega_\varepsilon$ is pluriclosed.
\end{theorem}
\begin{proof}
With the torus coordinates above set $s=\sin(2\pi x_1)$ and
\[
 \xi=\frac i2(dz_1\wedge d\bar z_2+dz_2\wedge d\bar z_1)
      =dx_1\wedge dy_2+dx_2\wedge dy_1,
 \qquad h_\tau=\beta+\tau s\xi.
\]
The first relative Hermitian block is
$\left(\begin{smallmatrix}1&\tau s\\\tau s&1\end{smallmatrix}\right)$,
so positivity follows from $|\tau s|<1$. Moreover
\[
 dh_\tau=2\pi\tau\cos(2\pi x_1)
                   dx_1\wedge dx_2\wedge dy_1\ne0.
\]
Because $\ddc s$ is proportional to $dz_1\wedge d\bar z_1$,
$\ddc(s\xi)=0$. The first two-coordinate metric is thus pluriclosed,
and its square is top degree on that complex surface. Expanding with
the remaining closed curve factors proves
\[
 \ddc(h_\tau^j)=0\quad(1\leq j\leq n-1),\qquad
 \int_X h_\tau^n=n!\int_0^1(1-\tau^2\sin^2(2\pi x_1))dx_1
                   =V_\tau.
\]
Use the same $u_\varepsilon$ and $a_\varepsilon$ as in
Theorem~\ref{th:degeneration}, now with $\varepsilon\leq1/2$. Then
\[
 \Omega_\varepsilon=a_\varepsilon\beta_1+
       \sum_{j\geq2}\beta_j+\tau s\xi,
 \qquad a_\varepsilon=\varepsilon+2(1-\varepsilon)s^2.
\]
The first block of $\Omega_\varepsilon-\varepsilon\beta$ has
diagonal entries $2(1-\varepsilon)s^2$ and $1-\varepsilon$, and
determinant
\[
 \bigl[2(1-\varepsilon)^2-\tau^2\bigr]s^2\geq0.
\]
It is positive semidefinite; the remaining diagonal entries are
$1-\varepsilon$. Hence $\Omega_\varepsilon\geq\varepsilon\beta$,
with equality in the least eigenvalue at $s=0$.
Also $\Omega_\varepsilon\leq3\beta$, all fixed-order smooth norms
are uniformly bounded, and the smooth semipositive limit loses rank
at $s=0$. Since $d\ddc u_\varepsilon=0$, the exterior derivative
is the displayed fixed nonzero $dh_\tau$.

Lemma~\ref{lem:mass} gives the lower-power identities for every
$\Omega_\varepsilon$, in particular pluriclosedness, and gives both
bounded-potential masses as $V_\tau$.
Theorem~\ref{th:profile} applied to $h_\tau$ gives equality of the
entire current and measure sets and their exact endpoint for every
$p\in L^n$. The balanced trace is fixed as well:
\[
 \int_X\Omega_\varepsilon\wedge\beta^{n-1}=n!.
\]
The off-diagonal term has zero trace and the mean of
$a_\varepsilon$ equals one. The global holomorphic volume form gives
$c_1=0$. The manifold itself is already K\"ahler, so these genuinely
nonclosed metrics furnish an estimate obstruction on a coexisting
manifold, rather than a counterexample to coexistence.
\end{proof}

\section{Concentration at sharp all-current maximizers}
Sharp determinant bounds also need not give uniform integrability of
the determinant powers or preserve maximizers under weak current limits.
\begin{theorem}\label{th:concentration}
On the fixed torus with $\beta,V,dV$ used above, there are smooth
positive closed $A_\varepsilon=\beta+\ddc w_\varepsilon$,
$0<\varepsilon<1$, and $T_\varepsilon=A_\varepsilon^{n-1}/V
\in\Ga_\beta$ satisfying
\[
 \int_X k_{T_\varepsilon}^{n/(n-1)}dV=V^{-1/(n-1)}.
\]
Each $T_\varepsilon$ is an actual dual maximizer for a specified
unit $L^n$ root. The potentials $w_\varepsilon$ have uniformly
bounded oscillation and converge uniformly to a continuous canonical
$\beta$-psh $w_0$. The currents converge weakly to
$T_0=(\beta+\ddc w_0)^{n-1}/V\in\Ga_\beta$, but
$\Gamma_0(T_0)=0$. The power densities above are not uniformly
integrable. Nevertheless every $w_\varepsilon$, including $w_0$,
has bounded-potential Monge--Amp\`ere mass $V$.
\end{theorem}
\begin{proof}
Put $r=1-\varepsilon$ and
\[
 P_r(x)=\frac{1-r^2}{1-2r\cos(2\pi x)+r^2},\qquad
 f_\varepsilon=\varepsilon+(1-\varepsilon)P_r.
\]
The positive periodic Poisson kernel has mean one and Fourier series
\[
 P_r=1+2\sum_{m\geq1}r^m\cos(2\pi mx).
\]
Define
\[
 w_\varepsilon=-\frac{2(1-\varepsilon)}{\pi}
       \sum_{m\geq1}\frac{r^m\cos(2\pi mx_1)}{m^2}.
\]
For fixed $\varepsilon>0$ this is smooth. Differentiation gives
$w_\varepsilon''=4\pi(f_\varepsilon-1)$, and hence
\[
 A_\varepsilon=f_\varepsilon(x_1)\beta_1+
        \sum_{j=2}^n\beta_j=\beta+\ddc w_\varepsilon>0.
\]
It is closed, and $\int A_\varepsilon^n=V$.
Theorem~\ref{th:profile} and the smooth-root formula give
\[
 k_{T_\varepsilon}=V^{-1/n}f_\varepsilon^{(n-1)/n},
 \qquad k_{T_\varepsilon}^{n/(n-1)}
       =V^{-1/(n-1)}f_\varepsilon.
\]
Thus the displayed power integral attains the upper bound over
\emph{all} currents, not a bound for a selected family. More
precisely, $p_\varepsilon=f_\varepsilon^{1/n}$ has $L^n$ norm one,
and $\int p_\varepsilon k_{T_\varepsilon}dV=V^{-1/n}$, so
$T_\varepsilon$ attains the original entire-family dual supremum
for this varying root. No common fixed density is asserted.

The Fourier series bounds $|w_\varepsilon|\leq\pi/3$ and implies
uniform convergence to
\[
 w_0=-\frac2\pi\sum_{m\geq1}\frac{\cos(2\pi mx_1)}{m^2}.
\]
This limit is continuous and canonically $\beta$-psh by distributional
positivity of $A_0=\beta+\ddc w_0$.
The probability measures $P_r(x_1)dx_1$ converge weakly to the
periodic Dirac mass at $x_1=0$, as follows from their Fourier
coefficients $r^{|m|}$. Thus
\[
 A_0=\delta_{x_1=0}\beta_1+\sum_{j=2}^n\beta_j,
\]
where the Dirac measure is tensored with Lebesgue measure in the
remaining real coordinates. Its bounded local products are genuine
Bedford--Taylor products. Locally the first-coordinate current
$\beta_1+\ddc w_0$ is pulled back from one complex coordinate,
so its powers of degree at least two vanish in the bounded-product
recursion. The remaining factors are smooth and closed.
Multilinear expansion therefore contains at most one first-coordinate
factor; the same expansion for $A_\varepsilon$ passes directly to
the distributional limit. Consequently
$T_\varepsilon\to T_0=A_0^{n-1}/V$ weakly. This current is
positive, closed and normalized by $\int\beta\wedge T_0=1$.

In its positive coefficient matrix the absolutely continuous part is
proportional to the first diagonal matrix below; the singular part
on $x_1=0$ is proportional to the second:
\[
 \operatorname{diag}(1,0,\ldots,0),\qquad
 \operatorname{diag}(0,1,\ldots,1).
\]
Their ranks are respectively
one and $n-1$, both less than $n$. With a measure dominating these
mutually singular parts, formula~\eqref{eq:intrinsic} therefore gives
$\Gamma_0(T_0)=0$. This is not the Monge--Amp\`ere measure of
$w_0$, which has total mass $V$ by the closed bounded recursion.

Away from $x_1=0$, $f_\varepsilon\to0$ pointwise; its integral
against $dV$ remains one. Thus the positive power densities are not
uniformly integrable: almost-everywhere convergence to zero and
uniform integrability would force their fixed positive integrals to
vanish. In particular $T\mapsto\int k_T^{n/(n-1)}dV$ is not
continuous under these weak current limits, even with uniformly
convergent continuous canonical potentials. The intrinsic measures
themselves satisfy $\Gamma_0(T_\varepsilon)\to0$ in total mass;
this follows from their uniform $L^{n/(n-1)}$ bound by splitting
off arbitrarily small neighborhoods of $x_1=0$. The discontinuity
concerns their \emph{power integrals}, not that last measure limit.
\end{proof}

\begin{remark}
Even the entire normalized Gauduchon-current determinant data,
together with a fixed positive lower and upper bounded-potential mass,
$c_1=0$ and uniform smooth coefficient bounds,
cannot alone force a uniform positive pointwise lower metric bound.
The metrics in Theorem~\ref{th:degeneration} have identical such data but a vanishing least
eigenvalue. Their volume densities are $Va_\varepsilon$, whose
positive lower bound also degenerates. Theorem~\ref{th:nonclosed}
extends the obstruction to genuinely nonclosed pluriclosed metrics
with fixed nonzero exterior derivative and a separate balanced metric.
The lower-power $dd^c$-closed hypothesis was checked directly for that
family; no closed-background Stokes formula was imposed on it.
Estimates involving an
actual flow, its initial pointwise positivity, or additional density
and derivative control are not excluded. No pluriclosed-flow
trajectory is constructed; the original coexistence and convergence
conjectures remain unresolved. No literature priority is claimed.
\end{remark}
\begin{thebibliography}{9}
\bibitem{BT} E. Bedford and B. A. Taylor,
\emph{A new capacity for plurisubharmonic functions}, Acta Math.
149 (1982), 1--40, equation (2.1), Theorem 2.1 and Corollary 2.2.
\bibitem{TW} V. Tosatti and B. Weinkove,
\emph{The complex Monge--Amp\`ere equation on compact Hermitian
manifolds}, \href{https://arxiv.org/abs/0910.1390v2}{arXiv:0910.1390v2},
Corollary 1.
\end{thebibliography}
\end{document}
